\documentclass[10pt,a4paper]{amsart}
\usepackage[margin=19mm]{geometry}
\usepackage{amsmath,amssymb,mathtools}
\usepackage[scr=rsfs,cal=euler]{mathalfa}
\usepackage{xfrac}
\usepackage[unicode,hidelinks,bookmarks=false]{hyperref}
\newcommand{\ie}{i.e.\ }
\newcommand{\cf}{cf.\ }
\newcommand{\resp}{respectively\ }
\newcommand{\Subsection}{\subsection}
\providecommand{\nobreakdash}{\nobreak\hbox{-}}
\setlength{\emergencystretch}{2em}
\multlinegap0pt
\usepackage{amssymb}
\newtheorem{prop}{Proposition}
\newtheorem{thm}{Theorem}
\title{Generalized regular representations of big wreath products}
\author{Eugene Strahov}
\date{Source: arXiv:2303.11671v2 (CC BY 4.0)}
\begin{document}
\maketitle

\noindent\emph{Source and licence.} Generalized regular representations of big wreath products. Version arXiv:2303.11671v2 (CC BY 4.0), distributed by arXiv under the Creative Commons Attribution 4.0 International licence, http://creativecommons.org/licenses/by/4.0/, as recorded in the arXiv licence metadata for this exact version and shown on the abstract page https://arxiv.org/abs/2303.11671v2. Full licence notice: \texttt{licenses/arxiv-2303.11671-CC-BY-4.0.txt}.\par\smallskip

\noindent\textit{Editorial scope.} Counted unit: the article's thm THEOREMCocycles (source0.tex lines 897--903). The article's complete original proof of it is delivered with it (source0.tex lines 904--1136, one \texttt{proof} environment of 233 lines). No mathematics is written, repaired or re-derived: the statement and the proof are the article's own lines, in the article's own notation, and no retained line is substituted or re-flowed.  Retained uncounted as the source-local context the proof needs: lines 528--533; lines 581--581.  Nothing was omitted from any retained span, so the package carries no omission record.  No retained line contains a citation, so the wrapper carries no bibliography and no bibliography entry.  No retained line contains a figure environment, an image-inclusion command or drawing code, so no image is delivered and the package declares no asset dependency.  Editorial formatting only, in addition: an AMS \texttt{amsart} preamble that restates the article's own declarations and macros used by the retained lines, namely the environments \texttt{prop}, \texttt{thm} on separate counters, together with the standard packages those lines need.  The retained environments print in this document as Proposition 1, Theorem 1 in order of appearance, whereas the article prints the same items with the numbering of its own full document.  The complete native source of the article as distributed in the arXiv e-print for this exact version is bundled unchanged as \texttt{sources/138969/source0.tex}.
\begin{prop}\label{PROPOSITIONEQUIVARIANCEPROJECTION}The projection $p_{n,n+1}: G\sim S(n+1)\longrightarrow G\sim S(n)$ is equivariant with respect to the two-sided action of $G\sim S(n)$, i.e.
$$
p_{n,n+1}\left((\kappa,\pi)(g,s)(h,t)\right)=(\kappa,\pi)p_{n,n+1}\left((g,s)\right)(h,t),
$$
for each $(g,s)\in G\sim S(n+1)$, and each $(\kappa,\pi)\in G\sim S(n)$, $(h,t)\in G\sim S(n)$.
\end{prop}

\subsection{The wreath product $G_{\infty}=G\sim S(\infty)$}\label{SubSectionWreathInfty}

\begin{thm}\label{THEOREMCocycles} There exist $k$ integer functions $C_{l}\left(x,W\right)$ on $\mathfrak{S}_{G}\times\left(G_{\infty}\times G_{\infty}\right)$ uniquely defined by the following property:
if $n$ is large enough so that $W\in G_n\times G_n$, then
\begin{equation}\label{Cocycles}
C_l\left(x,W\right)=\left[p_n\left(xW\right)\right]_{c_l}-\left[p_n(x)\right]_{c_l}.
\end{equation}
Here $l=1,\ldots,k$;  $p_n :\;\mathfrak{S}_G\longrightarrow G\sim S(n)$ is the natural projection, and the action of $G_{\infty}\times G_{\infty}$ on $\mathfrak{S}_{G}$ is defined as in Section \ref{SubSectionWreathInfty}.
\end{thm}

\begin{proof}We need to prove that the quantities $C_l(x,W)$ defined by equation (\ref{Cocycles}) do not depend on $n$ provided that $n$ is so large that
the element $W$ of $G_{\infty}\times G_{\infty}$ already belongs to $G_n\times G_n$.



Let $W$ be an element of $G_n\times G_n$. By Proposition \ref{PROPOSITIONEQUIVARIANCEPROJECTION} the projection $p_n$ is equivariant with respect to two-sided action of $G_n$. Thus in order to
prove Theorem \ref{THEOREMCocycles} it is enough to show that the condition
\begin{equation}\label{CocCondition}
\left[x_nW\right]_{c_l}-\left[x_n\right]_{c_l}=\left[x_{n+1}W\right]_{c_l}-\left[x_{n+1}\right]_{c_l}
\end{equation}
is satisfied for all $x_n\in G_n$, and for all $x_{n+1}\in G_{n+1}$ such that $x_n=p_{n,n+1}\left(x_{n+1}\right)$.
\begin{prop}\label{PROPOSITIONCoW}Assume that (\ref{CocCondition}) is satisfied for $W_1,\ldots, W_m\in G_n\times G_n$. Then
(\ref{CocCondition}) is also satisfied for the product $W_1\ldots W_m$. In other words, if
\begin{equation}\label{CocCondition1}
\left[x_nW_p\right]_{c_l}-\left[x_n\right]_{c_l}=\left[x_{n+1}W_p\right]_{c_l}-\left[x_{n+1}\right]_{c_l}
\end{equation}
holds true for each $p=1,\ldots,m$,  and for all $x_n\in G_n$, $x_{n+1}\in G_{n+1}$ such that $x_n=p_{n,n+1}\left(x_{n+1}\right)$, then
\begin{equation}\label{CocCondition2}
\left[x_nW_1\ldots W_m\right]_{c_l}-\left[x_n\right]_{c_l}=\left[x_{n+1}W_1\ldots W_m\right]_{c_l}-\left[x_{n+1}\right]_{c_l}
\end{equation}
holds true as well.
\end{prop}
\begin{proof} The proof is by induction. If $m=1$, then (\ref{CocCondition2}) turns into
(\ref{CocCondition1}) which holds true by our assumption in the statement of Proposition \ref{PROPOSITIONCoW}.
Assume that $m\geq 2$, and that
\begin{equation}\label{CocCondition3}
\left[x_nW_1\ldots W_{m-1}\right]_{c_l}-\left[x_n\right]_{c_l}=\left[x_{n+1}W_1\ldots W_{m-1}\right]_{c_l}-\left[x_{n+1}\right]_{c_l}
\end{equation}
is satisfied for  all $x_n\in G_n$, and for all $x_{n+1}\in G_{n+1}$ such that $x_n=p_{n,n+1}\left(x_{n+1}\right)$.
Denote
$$
\tilde{x}_n=x_nW_1\ldots W_{m-1},
$$
and observe that
\begin{equation}\label{P3211}
p_{n,n+1}\left(x_{n+1}W_1\ldots W_{m-1}\right)=p_{n,n+1}\left(x_{n+1}\right)W_1\ldots W_{m-1}=\tilde{x}_n,
\end{equation}
where we have used the equivariance of the projection $p_{n,n+1}$, see Proposition \ref{PROPOSITIONEQUIVARIANCEPROJECTION}. Now we write
\begin{equation}\label{P3212}
\left[x_nW_1\ldots W_{m}\right]_{c_l}-\left[x_n\right]_{c_l}=\left[\tilde{x}_nW_m\right]_{c_l}-\left[\tilde{x}_n\right]_{c_l}+\left[\tilde{x}_n\right]_{c_l}-\left[x_n\right]_{c_l}.
\end{equation}
By our assumption in the statement of  Proposition \ref{PROPOSITIONCoW}, and by (\ref{P3211}) the first difference in the right-hand side of equation (\ref{P3212}) can be rewritten as
\begin{equation}\label{P3213}
\left[\tilde{x}_nW_m\right]_{c_l}-\left[\tilde{x}_n\right]_{c_l}=\left[x_{n+1}W_1\ldots W_m\right]_{c_l}-\left[x_{n+1}W_1\ldots W_{m-1}\right]_{c_l}.
\end{equation}
The second difference in the right-hand side of equation (\ref{P3212}) is
\begin{equation}\label{P3214}
\left[\tilde{x}_n\right]_{c_l}-\left[x_n\right]_{c_l}=\left[x_nW_1\ldots W_{m-1}\right]_{c_l}-\left[x_n\right]_{c_l}.
\end{equation}
By assumption (\ref{CocCondition3}), it can be replaced by $\left[x_{n+1}W_1\ldots W_{m-1}\right]_{c_l}-\left[x_{n+1}\right]_{c_l}$. Then we conclude that
$$
\left[x_nW_1\ldots W_{m}\right]_{c_l}-\left[x_n\right]_{c_l}=\left[x_{n+1}W_1\ldots W_{m}\right]_{c_l}-\left[x_{n+1}\right]_{c_l}
$$
holds true as well. Proposition \ref{PROPOSITIONCoW} is proved.
\end{proof}
Proposition \ref{PROPOSITIONCoW} implies that it suffices to prove (\ref{CocCondition})
for $W$ of the form $\left(w,e_{G_n}\right)$, and for $W$ of the form $\left(e_{G_n},w\right)$
where $w\in G_n=G\sim S(n)$, and $e_{G_n}$ is the unit element of $G_n=G\sim S(n)$. Below we consider the second case. In this case we need to prove that
\begin{equation}
\left[wx_n\right]_{c_l}-\left[x_n\right]_{c_l}=\left[wx_{n+1}\right]_{c_l}-\left[x_{n+1}\right]_{c_l}
\end{equation}
for any $w=\left(\left(g_1,\ldots,g_n\right),s\right)\in G\sim S(n)$, any $x_n\in G\sim S(n)$, and any $x_{n+1}\in G\sim S(n+1)$ such that $p_{n,n+1}\left(x_{n+1}\right)=x_n$.
\begin{prop}\label{PROPOSITIONCoW1} With $x_n$ and $x_{n+1}$ as above,
\begin{equation}\label{ge}
\left[\left(\left(g_1,\ldots,g_n\right),e_{S(n)}\right)x_n\right]_{c_l}-\left[x_n\right]_{c_l}
=\left[\left(\left(g_1,\ldots,g_n\right),e_{S(n)}\right)x_{n+1}\right]_{c_l}-\left[x_{n+1}\right]_{c_l},
\end{equation}
where $e_{S(n)}$ denotes the unit element of $S(n)$, and $g_1$, $\ldots$, $g_n$ are arbitrary elements of $G$.
\end{prop}
\begin{proof} Assume that $x_n=(h,t)$, $x_{n+1}=\left(\tilde{h},t'\right)$, and $t'$ is obtained from $t$ by inserting $n+1$ into the existing cycle:
$$
i_1\rightarrow \ldots\rightarrow i_p\rightarrow i_{p+1}\rightarrow\ldots\rightarrow i_1\Longrightarrow
i_1\rightarrow\ldots\rightarrow i_{p}\rightarrow n+1\rightarrow i_{p+1}\rightarrow\ldots\rightarrow i_1.
$$
If $h=\left(h_1,\ldots,h_n\right)$,  and $n+1$ is inserted between the numbers $i_p$ and $i_{p+1}$,
then $\tilde{h}=\left(h_1,\ldots,h_{i_{p+1}}h_{n+1}^{-1},\ldots,h_n,h_{n+1}\right)$,
i.e. $\tilde{h}$ is obtained from $h$ by adding an additional element of $G$ to the list $\left(h_1,\ldots,h_n\right)$
(this additional element is denoted by $h_{n+1}$), and by multiplication of $h_{i_{p+1}}$ by $h_{n+1}^{-1}$ from the right. Note that  the cycle-product of $x_{n+1}$ corresponding to the cycle
\begin{equation}\label{cycleiinserted}
i_1\rightarrow i_2\rightarrow\ldots\rightarrow i_{p}\rightarrow n+1\rightarrow i_{p+1}\rightarrow\ldots\rightarrow i_1
\end{equation}
is $h_{i_1}\ldots\left(h_{i_{p+1}}h_{n+1}^{-1}\right)h_{n+1}h_{i_p}\ldots h_{i_2}$,
which is the same as the cycle-product of $x_n$ corresponding to the cycle
\begin{equation}\label{cyclei}
i_1\rightarrow \ldots\rightarrow i_p\rightarrow i_{p+1}\rightarrow\ldots\rightarrow i_1
\end{equation}
We conclude that
$$
\left[x_n\right]_{c_l}=\left[x_{n+1}\right]_{c_l},\;\; l=1,\ldots, k
$$
in this case.

Now we have
\begin{equation}\label{4.13.a}
\left(\left(g_1,\ldots,g_n\right),e_{S(n)}\right)x_n=\left(\left(g_1h_1,\ldots,g_nh_n\right),t\right),
\end{equation}
and
$$
\left(\left(g_1,\ldots,g_n\right),e_{S(n)}\right)x_{n+1}
=\left(\left(g_1h_1,\ldots,g_{i_{p+1}}h_{i_{p+1}}h_{n+1}^{-1},\ldots,
g_nh_n,h_{n+1}\right),t'\right).
$$
We check that the cycle-product of $\left(\left(g_1,\ldots,g_n\right),e_{S(n)}\right)x_n$ corresponding to the
cycle (\ref{cyclei}) is the same as that of $\left(\left(g_1,\ldots,g_n\right),e_{S(n)}\right)x_{n+1}$
corresponding to the cycle (\ref{cycleiinserted}). Therefore,
$$
\left[\left(\left(g_1,\ldots,g_n\right),e_{S(n)}\right)x_n\right]_{c_l}
=\left[\left(\left(g_1,\ldots,g_n\right),e_{S(n)}\right)x_{n+1}\right]_{c_l}, \;\; l=1,\ldots, k.
$$
Thus we conclude that if $t'$ is obtained from $t$ by inserting $n+1$ into an existing cycle of $t$, then
equation  (\ref{ge}) is satisfied.

Now assume that $t'$  is obtained from $t$ by adding a new cycle of the form $(n+1)$. Then
$\tilde{h}=\left(h_1,\ldots,h_n,h_{n+1}\right)$, and we have
$$
\left(\left(g_1,\ldots,g_n\right),e_{S(n)}\right)x_{n+1}
=\left(\left(g_1h_1,\ldots,g_nh_n,h_{n+1}\right),t'\right).
$$
As for the product $\left(\left(g_1,\ldots,g_n\right)e_{S(n)}\right)x_n$, it is given by equation (\ref{4.13.a}).
We find
$$
\left[x_{n+1}\right]_{c_l}-\left[x_n\right]_{c_l}=
 \left\{
 \begin{array}{ll}
 1, & h_{n+1}\in c_l, \\
 0, & \mbox{otherwise},
\end{array}
\right.
$$
and
$$
\left[\left(\left(g_1,\ldots,g_n\right),e_{S(n)}\right)x_{n+1}\right]_{c_l}
-\left[\left(\left(g_1,\ldots,g_n\right),e_{S(n)}\right)x_n\right]_{c_l}=
 \left\{
 \begin{array}{ll}
 1, & h_{n+1}\in c_l, \\
 0, & \mbox{otherwise}.
\end{array}
\right.
$$
Equation (\ref{ge}) holds true in this case as well.
\end{proof}
\begin{prop}\label{PROPOSITIONCoW2} Assume that $s\in S(n)$ is a transposition $(ij)$ (where $1\leq i<j\leq n$), and let $e_{G^{n}}$ denote the unit element of $G^n=G\times\ldots\times G$. We have
\begin{equation}\label{3.3.13}
\left[\left(e_{G^n},s\right)x_n\right]_{c_l}-\left[x_n\right]_{c_l}
=\left[\left(e_{G^n},s\right)x_{n+1}\right]_{c_l}-\left[x_{n+1}\right]_{c_l}
\end{equation}
for each $l=1,\ldots,k$, each $x_n\in G\sim S(n)$, and each $x_{n+1}\in G\sim S(n+1)$ such that $p_{n,n+1}\left(x_{n+1}\right)=x_n$.
\end{prop}
\begin{proof}Set
$$
x_n=\left(g,t\right),\; g=\left(g_1,\ldots,g_n\right),\; t\in S(n),
$$
and write $t$ as a product of cycles. First, let us assume that $i$ and $j$ are situated in two different cycles of $t$.
We can write these cycles explicitly as
\begin{equation}\label{3.3.14}
i\rightarrow i_1\rightarrow\ldots\rightarrow i_p\rightarrow i,
\end{equation}
and
\begin{equation}\label{3.3.15}
j\rightarrow j_1\rightarrow\ldots\rightarrow j_f\rightarrow j.
\end{equation}
The cycle-product of $x_n$ corresponding to the first cycle is $g_{i_p}\ldots g_{i_1}g_i$, and the cycle-product of $x_n$ corresponding to the second cycle is $g_{j_f}\ldots g_{j_1}g_{j}$.

Set $x_{n+1}=\left(g',t'\right)$, and assume that $t'$ is obtained from $t$ by creating a new cycle $(n+1)$, or by inserting $n+1$ to an existing cycle of $t$ which do not contain both $i$ and $j$. In this case it is not hard to check that (\ref{3.3.13}) is satisfied.
Indeed, the cycle in which $n+1$ is going to be inserted does not affect the difference
$\left[\left(e_{G^n},s\right)x_n\right]_{c_l}-\left[x_n\right]_{c_l}$, and the same cycle with inserted $n+1$ does not affect the difference
$\left[\left(e_{G^n},s\right)x_{n+1}\right]_{c_l}-\left[x_n\right]_{c_l}$.

If $t'$ is obtained from $t$ by inserting $n+1$ into the cycle with $i$, then  the cycle (\ref{3.3.14})
turns into the cycle
\begin{equation}\label{3.3.17}
i\rightarrow i_1\rightarrow\ldots\rightarrow i_m\rightarrow n+1\rightarrow i_{m+1}\rightarrow
\ldots\rightarrow i_p\rightarrow i,
\end{equation}
and we obtain $\left[x_{n+1}\right]_{c_l}=\left[x_n\right]_{c_l}$.  Indeed,
as soon as $x_n=p_{n,n+1}\left(x_{n+1}\right)$,
the cycle-product of $x_n$ corresponding to the cycle (\ref{3.3.14}) is the same as that of $x_{n+1}$ corresponding to the cycle
 (\ref{3.3.17}).

In $\left(e_{G^n},s\right)x_n$ the cycles (\ref{3.3.14}) and (\ref{3.3.15}) of $t$ merge into the cycle
\begin{equation}\label{3.3.18}
i\rightarrow i_1\rightarrow\ldots\rightarrow i_p\rightarrow j\rightarrow j_{1}\rightarrow
\ldots\rightarrow j_f\rightarrow i.
\end{equation}
The cycle-product of $\left(e_{G^n},s\right)x_n$ corresponding to the cycle (\ref{3.3.18})
is
\begin{equation}\label{zctype}
g_ig_{j_f}\ldots g_{j_1}g_{j}g_{i_p}\ldots g_{i_1}.
\end{equation}
In $\left(e_{G^n},s\right)x_{n+1}$ the cycles (\ref{3.3.17}) and (\ref{3.3.15}) of $t'$ merge into the cycle
\begin{equation}\label{3.3.19}
i\rightarrow i_1\rightarrow\ldots\rightarrow i_m\rightarrow n+1\rightarrow i_{m+1}\rightarrow
\ldots\rightarrow i_p\rightarrow j\rightarrow j_1\rightarrow \ldots\rightarrow j_f\rightarrow i,
\end{equation}
and the cycle-product of $\left(e_{G^n},s\right)x_{n+1}$ corresponding to this cycle is given by
(\ref{zctype}) as well. We conclude that
\begin{equation}\label{Zvzbv}
\left[\left(e_{G^n},s\right)x_{n}\right]_{c_l}=\left[\left(e_{G^n},s\right)x_{n+1}\right]_{c_l},
\end{equation}
and that equation (\ref{3.3.13}) holds true provided $i$ and $j$ are situated in two different cycles of $t$.

Second, consider the case where $i$ and $j$ are situated in the same cycle of $t$. Let us write this cycle as
\begin{equation}\label{3.3.20}
i\rightarrow i_1\rightarrow\ldots\rightarrow i_p\rightarrow j\rightarrow j_{1}\rightarrow
\ldots\rightarrow j_f\rightarrow i.
\end{equation}
In $t'$ the corresponding cycle includes $n+1$, and it takes the form
\begin{equation}\label{3.3.21}
i\rightarrow i_1\rightarrow\ldots\rightarrow i_{m}\rightarrow n+1\rightarrow i_{m+1}\rightarrow\ldots
\rightarrow i_p\rightarrow j\rightarrow j_{1}\rightarrow
\ldots\rightarrow j_f\rightarrow i.
\end{equation}
Again, we have
\begin{equation}
\left[x_{n}\right]_{c_l}=\left[x_{n+1}\right]_{c_l},
\end{equation}
for all $l=1,\ldots,k$ as the cycle-product of $x_n$ corresponding to (\ref{3.3.20}) is the same as that of $x_{n+1}$ corresponding to (\ref{3.3.21}). Now, in $\left(e_{G^n},s\right)x_{n}$ the multiplication of $x_n$ by $\left(e_{G^n},s\right)$ leads to the splitting of (\ref{3.3.20}) into two cycles
\begin{equation}\label{3.3.22}
i\rightarrow i_1\rightarrow\ldots\rightarrow i_p\rightarrow i\;\;\; \mbox{and}\;\;\; j\rightarrow j_{1}\rightarrow
\ldots\rightarrow j_f\rightarrow j.
\end{equation}
Moreover, in $\left(e_{G^n},s\right)x_{n+1}$ the multiplication of $x_{n+1}$ by $\left(e_{G^n},s\right)$ leads to the splitting of (\ref{3.3.21}) into two cycles
\begin{equation}\label{3.3.23}
i\rightarrow i_1\rightarrow\ldots\rightarrow i_m\rightarrow n+1\rightarrow i_m\rightarrow\ldots\rightarrow i\;\;\; \mbox{and}\;\;\; j\rightarrow j_{1}\rightarrow
\ldots\rightarrow j_f\rightarrow j.
\end{equation}
Again we check that the cycle-products of   $\left(e_{G^n},s\right)x_n$ corresponding to (\ref{3.3.22}) are the same as those of $\left(e_{G^n},s\right)x_{n+1}$ corresponding to (\ref{3.3.23}), and conclude that
(\ref{Zvzbv}) is satisfied. Therefore, equation (3.3.13) holds true in the situation where $i$ and $j$ are situated in the same cycle of $t$ as well.
\end{proof}
Propositions (\ref{PROPOSITIONCoW}), (\ref{PROPOSITIONCoW1}), and (\ref{PROPOSITIONCoW2}) imply that equation (\ref{CocCondition}) is satisfied for $W=\left(e_{G_n},w\right)$, where $w$ is an arbitrary element of $G_n$. The case of  $W=\left(w,e_{G_n}\right)$ can be considered in the same way.
Theorem \ref{THEOREMCocycles} is proved.
\end{proof}

\end{document}
